\section{Static graph identities and differential estimates}
\label{sec:graph-calculus}

This Section proves the differential identities used in both constructions.
They are identities on a Riemannian base; no spacetime satisfying field
equations is assumed.  In particular, the reference tensor in this Section
may be either the given maximal tensor or the synthetic tensor constructed
later.

The curvature, current, and boundary identities give the scalar, mass,
and boundary inequalities in both deformations. The maximum estimate
controls finite slopes without differentiating the trace cutoff; the
energy and Newton-current estimates supply the seam-continuity argument
in Section~\ref{sec:transmission-analysis}.

Let $(U,g)$ be a smooth three-dimensional Riemannian manifold.  Unless a
metric is displayed, covariant derivatives, norms, contractions, volume
forms, and divergences in this Section use $g$.  For a symmetric covariant
tensor $C$, the notation $C p$ means the vector obtained by contracting with
$p$ and raising the remaining index with $g$.  For a positive smooth
function $N$ and a smooth function $f$, define $s>0$ and $p$ by
\begin{equation}
 p=s^2\nabla f,\qquad s^2+|\dd f|^2s^4=N^2,
 \qquad v=\frac pN,\qquad W=\frac Ns,
 \qquad \bar g=g+s^2\dd f\otimes\dd f.
 \label{eq:calculus-constitutive}
\end{equation}
The positive root is unique, and it depends smoothly on $(N,\dd f)$,
including at $\dd f=0$, because
\begin{equation}
 s^2=\frac{2N^2}{1+\sqrt{1+4N^2|\dd f|^2}}.
 \label{eq:calculus-positive-root}
\end{equation}
We use $d=s/N=\sqrt{1-|v|^2}$ for a scalar in $(0,1]$; differentials are
written $\dd$.  Direct rank-one inversion and the determinant formula give
\begin{equation}
 \bar g^{-1}=g^{-1}-v\otimes v,\qquad
 \dd V_{\bar g}=W\,\dd V_g.
 \label{eq:calculus-volume-inverse}
\end{equation}
For a further positive function $w$ put
\begin{equation}
 x=sw,\qquad h=x\bar g,\qquad A=Ns\bar g^{-1},
 \qquad \pi=N^{-1}\operatorname{sym}\nabla p^\flat,
 \label{eq:calculus-conformal-definitions}
\end{equation}
where $\operatorname{sym}$ is the half-sum.  The tensor $\pi$ throughout
this Section is the graph second fundamental form, not the electrostatic
stress introduced in Section~\ref{sec:stress-extension}.  The elementary
identities
\begin{equation}
 p\cdot\dd f=W^2-1,\qquad
 A(\dd f,\dd f)=W-W^{-1}\le p\cdot\dd f,
 \qquad 0<A\le N^2g^{-1}
 \label{eq:calculus-energy-algebra}
\end{equation}
will also be used in the penalty estimates.

\subsection{Curvature and the conserved current}

\begin{lemma}[Static graph curvature]\label{lem:graph-curvature}
Suppose that $\divg_g p=TN$, where $T$ is a smooth scalar function.
Then $\tr_g\pi=T$ and
\begin{align}
 \bar\nabla\log s&=\nabla\log N-\pi(v),
 \label{eq:calculus-lapse-differential}\\
 R_{\bar g}
 &=R_g+T^2-|\pi|^2-2(\divg_g\pi-\dd T)\cdot v.
 \label{eq:calculus-static-curvature}
\end{align}
More generally, let $\mathcal K$ be any smooth symmetric covariant tensor,
and define
\begin{equation}
 \kappa=\tr_g\mathcal K,\qquad
 \ell=T-\kappa,\qquad \lambda=\pi-\mathcal K,
 \qquad C=\lambda-\ell g,
 \label{eq:calculus-trace-definitions}
\end{equation}
and
\begin{equation}
 \mathscr D(v)=R_g+\kappa^2-|\mathcal K|^2
       -2(\divg_g\mathcal K-\dd\kappa)\cdot v.
 \label{eq:calculus-dominant-expression}
\end{equation}
Then
\begin{equation}
 R_{\bar g}=\mathscr D(v)+|\lambda|^2-\ell^2
                   -\frac2N\divg_g(Cp).
 \label{eq:calculus-completed-curvature}
\end{equation}
\end{lemma}

\begin{proof}
The trace assertion follows immediately from the definition of $\pi$.
To prove the first identity without any geometric interpretation, set
$a=\dd\log s$.  Since $p^\flat=s^2\dd f$,
\[
 \nabla_i p_j=N\pi_{ij}+a_i p_j-a_j p_i.
\]
Differentiating $s^2=N^2-|p|^2$ and substituting this formula gives
\[
 N^2a_i-(a\cdot p)p_i=N\nabla_iN-N\pi_{ij}p^j.
\]
Division by $N^2$ and Equation~\eqref{eq:calculus-volume-inverse} give
Equation~\eqref{eq:calculus-lapse-differential}, as an equality of vectors.

For the curvature identity consider, solely for this computation, the
Lorentzian metric
\[
 \widetilde g=-s^2\dd t^2+\bar g
\]
on $\R\times U$.  The graph $t=f$ has induced metric $g$.  If
$X^\mathrm{tan}=X+\dd f(X)\partial_t$, its future unit normal is
$n=(\partial_t+p^\mathrm{tan})/N$.  Thus
$\partial_t=Nn-p^\mathrm{tan}$.  The tangential part of the Killing
equation for $\partial_t$ is
$N\widetilde g(\widetilde\nabla_{X^\mathrm{tan}}n,Y^\mathrm{tan})
=\operatorname{sym}\nabla p^\flat(X,Y)$.
Its second fundamental form, with the sign convention specified after
\eqref{eq:future-trapping},
is therefore $\pi$.

Write $\widetilde G$ for the Einstein tensor and
$n_0=s^{-1}\partial_t$ for the normal to a static slice.  The contracted
Gauss and Codazzi identities with timelike unit normal and second form
$\widetilde g(\widetilde\nabla n,\cdot)$ read
\[
 2\widetilde G(n,n)=R_g+T^2-|\pi|^2,
 \qquad
 \widetilde G(n,X^\mathrm{tan})
       =(\divg_g\pi-\dd T)(X).
\]
The static slices have zero second form.  Applying these same contractions
to a static slice gives
$\widetilde G(n_0,n_0)=R_{\bar g}/2$ and
$\widetilde G(n_0,X)=0$ for spatial $X$.  The static time component of
$n$ is $N/s$, whereas that of $\partial_t/N$ is $s/N$.  Consequently
\[
 \frac12R_{\bar g}
 =\widetilde G(n,\partial_t/N)
 =\widetilde G(n,n)-\widetilde G(n,v^\mathrm{tan}),
\]
which proves Equation~\eqref{eq:calculus-static-curvature} with the
claimed sign of the momentum term.

For completeness, the completion of the square uses symmetry to give
\[
 \frac1N\divg_g(Cp)
 =(\divg_g\lambda-\dd\ell)\cdot v
                         +\lambda:\pi-\ell T.
\]
Also $2\ell T-\ell^2=T^2-\kappa^2$ and
$|\lambda|^2-2\lambda:\pi=|\mathcal K|^2-|\pi|^2$.
Substituting these two equalities in the right side of
Equation~\eqref{eq:calculus-completed-curvature} recovers
Equation~\eqref{eq:calculus-static-curvature} term by term.
\end{proof}

\begin{lemma}[Graph current and conformal curvature]\label{lem:graph-current}
Under the hypotheses and notation of Lemma~\ref{lem:graph-curvature}, set
\begin{equation}
 B=\nabla N-(\mathcal K+\ell g)p,\qquad
 \divg_g B=\Phi,\qquad -\divg_g(A\dd w)=\Psi,
 \qquad Q=B+A\dd w.
 \label{eq:calculus-general-system}
\end{equation}
Then
\begin{align}
 Q&=N\bar\nabla x+xCp-(x-1)B,\notag\\
 \divg_g Q&=\Phi-\Psi,
 \label{eq:calculus-general-current}\\
 N\Delta_{\bar g}x+x\divg_g(Cp)&=x\Phi-\Psi,
 \label{eq:calculus-general-x}\\
 xR_h&=\mathscr D(v)+|\lambda|^2-\ell^2
       -\frac{2\Phi}{N}+\frac{2\Psi}{Nx}
       +\frac32|\dd\log x|_{\bar g}^2.
 \label{eq:calculus-general-conformal}
\end{align}
In particular, the following two specializations hold.

\smallskip\noindent
\emph{The maximal AH system.}  If $\tr_gK=0$ and
\begin{equation}
 \divg_g p=0,\quad B=\nabla N-Kp,\quad
 \divg_g B=3N+\Sigma,\quad
 -\divg_g(A\dd w)=3N+\Sigma x+P,
 \label{eq:calculus-ah-system}
\end{equation}
then, writing
$D=R_g+6-|K|^2-2\divg_gK\cdot v$, one has
\begin{align}
 Q&=N\bar\nabla x+x(\pi-K)p-(x-1)B,\notag\\
 \divg_g Q&=-(x-1)\Sigma-P,
 \label{eq:calculus-ah-current}\\
 \left(\Delta_{\bar g}-3+
          N^{-1}\divg_g((\pi-K)p)\right)x&=-3-P/N,
 \label{eq:calculus-ah-x}\\
 x^2(R_h+6)
 &=x\bigl(D+|\pi-K|^2\bigr)
   +\frac{3}{2x}|\dd x|_{\bar g}^2
   +6(x-1)^2+\frac{2P}{N}.
 \label{eq:calculus-ah-curvature}
\end{align}
For the constraint densities in \eqref{eq:constraints},
$D=16\pi(\mu-J\cdot v)\ge0$ under the dominant energy condition.

\smallskip\noindent
\emph{The compact trace system.}  Let $0\le u_*\le1$ be constant and
take $\mathcal K=u_*K$, $\kappa=u_*\tr_gK$, $\ell=T-\kappa$.
If
\begin{equation}
 \begin{aligned}
 \divg_g p&=TN,&
 B&=\nabla N-(u_*K+\ell g)p,\\
 \divg_g B&=\Sigma_++\Sigma_-,&
 -\divg_g(A\dd w)&=\Sigma_+x+P,
 \end{aligned}
 \label{eq:calculus-compact-system}
\end{equation}
then
\begin{align}
 Q&=N\bar\nabla x+x(\lambda-\ell g)p-(x-1)B,\notag\\
 \divg_gQ&=-(x-1)\Sigma_+-P+\Sigma_-,
 \label{eq:calculus-compact-current}\\
 xR_h&=\mathscr D(v)+|\lambda|^2-\ell^2
             -\frac{2\Sigma_-}{N}+\frac{2P}{Nx}
             +\frac32|\dd\log x|_{\bar g}^2.
 \label{eq:calculus-compact-curvature}
\end{align}
Here $\lambda=\pi-u_*K$.  At $u_*=1$ this is the scalar identity for
the synthetic data.  It also holds on the outer region with
$p=T=K=\ell=0$, $\Sigma_-=0$, and $\mathscr D=R_{g_o}$.
\end{lemma}

\begin{proof}
Since $w=x/s$, Equation~\eqref{eq:calculus-lapse-differential} gives
\[
 A\dd w=N\bar\nabla x-Nx\bar\nabla\log s
        =N\bar\nabla x-x\nabla N+x\pi p.
\]
Adding $B$ proves the first formula for $Q$.  Its divergence from the
defining equations is $\Phi-\Psi$.  To compute the divergence of its
other expression, Equation~\eqref{eq:calculus-volume-inverse} gives
\begin{align*}
 \divg_g(N\bar\nabla x)
 &=N\Delta_{\bar g}x
       +N\bar g^{-1}(\dd\log s,\dd x)\\
 &=N\Delta_{\bar g}x+(\nabla N-\pi p)\cdot\dd x.
\end{align*}
The derivative-of-$x$ terms in $\divg_g(xCp+(1-x)B)$ are
$(Cp-B)\cdot\dd x=(\pi p-\nabla N)\cdot\dd x$.
They cancel exactly, leaving
\[
 \divg_gQ=N\Delta_{\bar g}x+x\divg_g(Cp)+(1-x)\Phi.
\]
This proves Equation~\eqref{eq:calculus-general-x}.  In dimension three
the conformal change $h=x\bar g$ has scalar curvature
\[
 xR_h=R_{\bar g}-2x^{-1}\Delta_{\bar g}x
                     +\frac32x^{-2}|\dd x|_{\bar g}^2.
\]
Substitution of Equations~\eqref{eq:calculus-completed-curvature} and
\eqref{eq:calculus-general-x} cancels the two occurrences of
$2N^{-1}\divg_g(Cp)$ and proves
Equation~\eqref{eq:calculus-general-conformal}.
The two stated systems now follow by substituting their values of
$T,\mathcal K,\Phi,\Psi$.  In the AH case, the nondifferential
cosmological terms after multiplication by $x$ are
$6x^2-12x+6=6(x-1)^2$.
Finally, $|v|<1$ and the constraint definitions give the asserted sign
of $D$.
\end{proof}

For later use, the divergence theorem makes these current formulas into
exact finite-domain flux budgets.  Let $S$ be the inner boundary, with
normal pointing into the domain, let $S_R$ have the outward normal, and
suppose $B_n=-b_S$ and $(A\dd w)_n=0$ on $S$.  The scalar $b_S$ is a
boundary datum, unrelated to the hyperbolic reference metric $b$.
In the compact construction assume also that the two currents have
matching normal flux at their common face, measured toward the end.
Then
\begin{equation}
 -\int_{S_R}Q_n\,\dd A=
 \begin{cases}
 \displaystyle\int_S b_S\,\dd A+
    \int_{U_R}\bigl[(x-1)\Sigma+P\bigr]\,\dd V,
       &\text{for \eqref{eq:calculus-ah-system}},\\[3pt]
 \displaystyle\int_S b_S\,\dd A+
    \int_{U_R}\bigl[(x-1)\Sigma_++P-\Sigma_-\bigr]\,\dd V,
       &\text{for \eqref{eq:calculus-compact-system}}.
 \end{cases}
 \label{eq:calculus-flux-budget}
\end{equation}
Indeed, the outward normal at $S$ is $-n$, so its contribution to the
outward flux is $\int_S b_S$.  The two face contributions cancel since
the base induced area forms and the directed fluxes agree.  Conversion
of the outer flux into the relevant mass difference requires the end
estimates proved in the corresponding construction.

\subsection{Boundary identities}

\begin{lemma}[Constant-height boundary]\label{lem:graph-boundary}
Let $\mathcal S$ be a smooth two-sided hypersurface on which $f$ is
constant.  Choose a $g$-unit normal $n$, and write $H=\divg_{\mathcal S}n$,
$\hat z=v\cdot n$, $y=\hat z^2$, and
$k_T=\tr_{\mathcal S}\mathcal K$.  The same directed normal is used on
each side if $\mathcal S$ is a gluing face.  Then
\begin{align}
 d^2\partial_n\log s&=B_n/N-k_T\hat z+yH,
 \label{eq:calculus-boundary-lapse}\\
 s\sqrt{x}\,H_h
 &=N(H-k_T\hat z)+B_n+(A\dd w)_n/x.
 \label{eq:calculus-boundary}
\end{align}
In particular, in the maximal AH system with $v=-z n$,
$k=K(n,n)$, $B_n=-b_S$, and $(A\dd w)_n=0$,
\begin{equation}
 s\sqrt{x}\,H_h=N(H-zk)-b_S=NH+\partial_nN.
 \label{eq:calculus-ah-boundary}
\end{equation}
In the compact system, if $H\le0$, $k_T<0$, $\hat z\le0$,
$B_n=-b_*<0$, and $(A\dd w)_n=0$, then $H_h<0$.

Suppose, at a compact gluing face, that $N_o=s_i$, that $w$ and the
normal fluxes of $B,A\dd w$ agree, that $f$ is constant on the inner
face, and that $p=0$ outside.  The induced $h$-metrics then agree.
If the inner base mean curvature satisfies $H_i>|k_T|$ and the outer
one is $H_o=\sqrt{H_i^2-k_T^2}$, then
$H_{h,i}\ge H_{h,o}$, both measured toward the end.
\end{lemma}

\begin{proof}
On $\mathcal S$, $p=N\hat z n$.  For tangent vectors $X,Y$,
$\nabla^2f(X,Y)=f_n\operatorname{II}(X,Y)$; consequently
$\pi(X,Y)=\hat z\operatorname{II}(X,Y)$ and
$\pi(n,n)=T-\hat z H$.  Taking the normal component in
Equation~\eqref{eq:calculus-lapse-differential} therefore yields
\[
 d^2\partial_n\log s
 =\partial_n\log N-\hat z T+yH.
\]
Since $\partial_nN=B_n+N(\mathcal K(n,n)+\ell)\hat z$ and
$\mathcal K(n,n)+\ell-T=-k_T$, this proves
Equation~\eqref{eq:calculus-boundary-lapse}.

In base normal coordinates, the tangential entries of $\bar g$ and
their first normal derivatives agree with those of $g$ at the face:
the extra products contain at least one tangential derivative of $f$.
The mixed entries vanish along the entire face, and the normal length
factor is $W$.  Hence $\bar n=dn$ and $H_{\bar g}=dH$.  The
mean-curvature formula for $h=x\bar g$ gives
\[
 \sqrt{x}\,H_h=d(H+\partial_n\log x).
\]
Multiplication by $s=Nd$, followed by
Equation~\eqref{eq:calculus-boundary-lapse}, gives
\[
 s\sqrt{x}\,H_h
 =N(H-k_T\hat z)+B_n+Nd^2w_n/w.
\]
The final term equals $(A\dd w)_n/x$, proving
Equation~\eqref{eq:calculus-boundary}.  In the maximal case
$k_T=-k$, $\hat z=-z$, and $\partial_nN=-Nzk-b_S$;
this proves Equation~\eqref{eq:calculus-ah-boundary}.
The asserted strict sign in the compact case follows because
$H-k_T\hat z\le0$ and $B_n<0$.

At the gluing face the tangential graph metrics equal their base
metrics, which agree, and $x_i=s_iw=N_ow=x_o$.
The remaining inequality reduces by
Equation~\eqref{eq:calculus-boundary} to
\[
 N_i(H_i-k_T\hat z)\ge s_i\sqrt{H_i^2-k_T^2}.
\]
Both sides are nonnegative, and after division by $N_i$ its squared
difference is
\[
 (H_i-k_T\hat z)^2-(1-\hat z^2)(H_i^2-k_T^2)
       =(k_T-H_i\hat z)^2\ge0.
\]
This proves the mean-curvature ordering.
\end{proof}

\subsection{A maximum estimate without differentiating the cutoff}

\begin{lemma}[Gradient maximum estimate]\label{lem:gradient-maximum}
Let $N>0,f,\mathcal K,T$ be smooth and satisfy
$\divg_g p=TN$ and
$\divg_g[\nabla N-(\mathcal K+\ell g)p]=\Phi$, with the definitions
above.  Suppose on the region in question that
\begin{equation}
 |\Ric_g|+|\mathcal K|+
       |\divg_g\mathcal K-\dd\tr_g\mathcal K|+|T|\le C_0,
 \qquad \Phi\le C_0N+\Sigma,\qquad \Sigma\ge0.
 \label{eq:calculus-maximum-hypotheses}
\end{equation}
At a point where $m=|\dd f|>0$, set
\begin{equation}
 \begin{aligned}
 l&=\log m,\quad u=\log N,\quad e=\nabla f/m,\quad y=|v|^2,\\
 q&=\frac{1-y}{1+y},\quad \beta=1+q,\quad
 \mathcal P=g^{-1}+(q-1)e\otimes e,\quad G=-\log s.
 \end{aligned}
 \label{eq:calculus-gradient-notation}
\end{equation}
Let $j$ be twice continuously differentiable with $j''>0$, with either
sign allowed for $j'$, and let $\eta$ be a twice continuously
differentiable localization function.  At an interior local maximum of
$G+j(f)+\eta$ with $y\ge1/2$,
\begin{equation}
 j''q m^2\le C\left[1+\frac\Sigma N
       +(1-y)^2(j'm)^2+|\dd\eta|^2+|\nabla^2\eta|\right],
 \label{eq:calculus-gradient-maximum}
\end{equation}
where $C$ depends only on $C_0$.  In particular, at that point,
\begin{equation}
 W^2\le\frac C{j''}\left[
     N^2+N\Sigma+(j')^2
       +N^2\bigl(|\dd\eta|^2+|\nabla^2\eta|\bigr)\right].
 \label{eq:calculus-gradient-weighted}
\end{equation}
No positive lower bound for $N$, no derivative bound for $T$, and no
second derivative bound for $N$ enter these two estimates.

At a constant-height boundary, with the notation of
Lemma~\ref{lem:graph-boundary}, the exact normal derivative is
\begin{equation}
 d^2\partial_n\bigl(G+j(f)+\eta\bigr)
   =-B_n/N+k_T\hat z-yH+j'\hat z/N+d^2\partial_n\eta.
 \label{eq:calculus-gradient-boundary}
\end{equation}
For the maximal AH inner boundary, if $H\le k$, $b_S\ge0$, and
$v=-zn$, this implies
\begin{equation}
 G_n=-\frac{yH+u_n}{1-y}
       =\frac{b_S/N+zk-z^2H}{1-z^2}
       \ge\frac{zk}{1+z}.
 \label{eq:calculus-ah-gradient-boundary}
\end{equation}
\end{lemma}

\begin{proof}
Use a $g$-orthonormal frame at the point with first vector $e$ and
transverse indices $a,b\in\{2,3\}$.  Put
$U_{ab}=\nabla^2_{ab}f/m$ and $C_T=T/z$, where $z=\sqrt y$.
Differentiation of the constitutive law gives
\begin{equation}
 \begin{aligned}
 Nm&=\frac z{1-y},\qquad
 q'=\frac{\dd q}{\dd(l+u)}=-\frac{4yq}{(1+y)^2},\\
 \mathcal P:\nabla^2f&=mC_T-\beta\nabla f\cdot\dd u,
 \qquad \tr U=C_T-ql_1-\beta u_1.
 \end{aligned}
 \label{eq:calculus-height-nondivergence}
\end{equation}
Here and below $q'$ differentiates the scalar constitutive function,
not a spatial coordinate.  The derivative of the height flux with
respect to $\dd f$ is $s^2\mathcal P$; its derivative with respect to
$N$ is $s^2(\beta/N)\nabla f$.  This also proves that the
nondivergence equation is smooth and elliptic through zero gradient.

Differentiating the height equation in direction $e$, commuting the
third derivatives of a scalar, and differentiating its norm gives
\begin{align}
 \mathcal P:\nabla^2l={}&\Ric(e,e)-\beta u_{11}
   -q'(l_1+u_1)^2-\beta\dd l\cdot\dd u
   +|U|^2+(1-q)|\dd_\perp l|^2-ql_1^2 \notag\\
 &+\frac{\partial_1T}{z}
        +C_T\bigl[(1-q)l_1-qu_1\bigr].
 \label{eq:calculus-log-gradient-equation}
\end{align}
Here is an explicit account of the norm and coefficient terms in this
calculation.  In $m^{-1}\mathcal P:\nabla^2m$, the second derivatives
of the norm contribute $q|\dd_\perp l|^2+|U|^2$.
The contraction $m^{-1}(\nabla_e\mathcal P):\nabla^2f$ is
$q'(l_1+u_1)l_1+2(q-1)|\dd_\perp l|^2$.
Differentiating $-\beta m u_1$ contributes
$-\beta u_{11}-\beta\dd l\cdot\dd u
-q'(l_1+u_1)u_1$.
Passing from $m$ to $l$ subtracts
$\mathcal P(\dd l,\dd l)=ql_1^2+|\dd_\perp l|^2$.
Finally,
$m^{-1}\partial_1(mT/z)=T_1/z+C_T[(1-q)l_1-qu_1]$.
These are precisely all terms of
Equation~\eqref{eq:calculus-log-gradient-equation}.

Set $I=(1+y)^{-1}$ and $J=y/(1+y)$.  Then
$G_l=J$, $G_u=-I$, and all four second partial derivatives of $G$
as a function of $(l,u)$ equal $-q'/2$.
Since $-I(q-1)=J\beta$, the Hessian terms in $u$ in
$\mathcal P:\nabla^2G$ reduce exactly to $-I\Delta u$.
The lapse equation, expanded using the symmetry of
$\mathcal K+\ell g$, is
\begin{equation}
 \Delta u=\Phi/N+
     [\divg_g\mathcal K-\dd\kappa+\dd T]\cdot v
     +(\mathcal K+\ell g):\pi-|\dd u|^2.
 \label{eq:calculus-log-lapse}
\end{equation}
The coefficient of $T_1$ is $J/z-Iz=0$.  Thus no derivative of $T$
remains.  More precisely, with $a_y=-q'/2$, the resulting identity is
\begin{align}
 \mathcal P:\nabla^2G={}&J\Ric(e,e)-I\Phi/N
   -I(\divg_g\mathcal K-\dd\kappa)\cdot v
   -I(\mathcal K+\ell g):\pi \notag\\
 &+\mathcal Q+JC_T[(1-q)l_1-qu_1],
 \label{eq:calculus-master-gradient}\\
 \mathcal Q={}&J|U|^2+a_y|\dd l+\dd u|^2
  -J\beta\dd l\cdot\dd u
  +J(1-q)|\dd_\perp l|^2-Jql_1^2+I|\dd u|^2.
 \label{eq:calculus-gradient-quadratic}
\end{align}
In obtaining the coefficient $a_y$ of the normal square, one uses
$2J+q=1$: the two contributions are $-Jq'$ and $-qq'/2$.

At the specified maximum put $D_0=j'm$.  The vanishing first derivative
is equivalent to
\begin{equation}
 J(l_i+u_i)=u_i-D_0\delta_{i1}-\eta_i.
 \label{eq:calculus-maximum-substitution}
\end{equation}
First take $\dd\eta=0$ and separate
$|U|^2=|U_{\rm tf}|^2+(\tr U)^2/2$.
The portion of
$\mathcal Q+j'\mathcal P:\nabla^2f$ independent of $C_T$ is exactly
\begin{align}
 J|U_{\rm tf}|^2
 &+\frac{(3-2y)u_1^2-6(1-y)D_0u_1
                         +3(1-y)^2D_0^2}{2y(1+y)} \notag\\
 &+\frac{2-y}{y(1+y)}|\dd_\perp u|^2.
 \label{eq:calculus-maximum-positive-form}
\end{align}
This follows by inserting
$l_1=(Iu_1-D_0)/J$, $l_a=Iu_a/J$, and
$\tr U=C_T+[-u_1+(1-y)D_0]/y$ in
Equation~\eqref{eq:calculus-gradient-quadratic}.
The terms containing $C_T$, including the last term of
Equation~\eqref{eq:calculus-master-gradient}, are
\begin{align}
 \frac J2 C_T^2+
 C_T[J(1-2q)l_1-J(\beta+q)u_1+D_0]
   =\frac J2C_T^2+qC_T(2D_0-u_1).
 \label{eq:calculus-trace-cancellation}
\end{align}
This last simplification is the gain needed for a bounded trace target:
$q|D_0|\le(1-y)|D_0|$, rather than an unweighted $|D_0|$.

For $y\in[1/2,1)$ the positive coefficients in
Equation~\eqref{eq:calculus-maximum-positive-form} are bounded below.
Young's inequality therefore bounds that expression below by
$c(|U_{\rm tf}|^2+|\dd u|^2)-C(1-y)^2D_0^2$.
The trace terms in Equation~\eqref{eq:calculus-trace-cancellation}
have the same lower bound after reducing $c$ and adding a constant,
because $|C_T|\le\sqrt2 C_0$.
For general $\dd\eta$, use $D_0+\eta_1$ in the normal substitutions.
The unchanged term $-\beta D_0u_1$ adds $\beta\eta_1u_1$ relative to
that substitution.  The transverse substitutions add terms bounded by
$C(|\dd\eta|^2+|\dd\eta||\dd u|)$, and the additional trace term is
$-(1-2q)C_T\eta_1$.  Another application of Young's inequality thus
adds only $C|\dd\eta|^2$ to the lower bound.

The components of the graph second form are
\begin{equation}
 \pi_{11}=T-z\tr U,\qquad \pi_{ab}=zU_{ab},\qquad
 \pi_{1a}=\frac{z\beta}{2}(u_a+l_a).
 \label{eq:calculus-second-form-components}
\end{equation}
Equation~\eqref{eq:calculus-maximum-substitution} also gives
\[
 |\tr U|\le C\bigl(1+|\dd u|+(1-y)|D_0|+|\dd\eta|\bigr).
\]
Thus the term $(\mathcal K+\ell g):\pi$ is absorbed in the same
positive quadratic expression, with an error
$C[1+(1-y)^2D_0^2+|\dd\eta|^2]$.
The curvature and momentum terms are bounded, and
$-I\Phi/N\ge-C-\Sigma/N$.
At a local maximum the positive definite tensor $\mathcal P$ has
nonpositive contraction with $\nabla^2(G+j(f)+\eta)$.
Its remaining height term is $j''\mathcal P(\dd f,\dd f)=j''qm^2$,
and $|\mathcal P:\nabla^2\eta|\le C|\nabla^2\eta|$.
This proves Equation~\eqref{eq:calculus-gradient-maximum}.

Multiply it by $N^2$ and use
$N^2qm^2=J/(1-y)$, $J\ge1/3$, and
$N^2(1-y)^2(j'm)^2=y(j')^2$.
This proves Equation~\eqref{eq:calculus-gradient-weighted} without a
lapse floor.  Lastly, negate
Equation~\eqref{eq:calculus-boundary-lapse} and use
$d^2f_n=\hat z/N$ to obtain
Equation~\eqref{eq:calculus-gradient-boundary}.
In the AH case, $\tr\pi=0$ also gives
$d^2\partial_n\log s=u_n+yH$.
Substitution of $u_n=-zk-b_S/N$ and $H\le k$ proves
Equation~\eqref{eq:calculus-ah-gradient-boundary}.
\end{proof}

\subsection{Energy and Newton flux identities at finite slope}

\begin{lemma}[Differential energy inequality]\label{lem:gradient-energy}
Assume the two differential equations and the geometric coefficient
bounds of Lemma~\ref{lem:gradient-maximum}.  On a region suppose, in
addition, that $N,s$ have positive lower bounds and finite upper bounds,
that $|\dd f|$ has a finite upper bound, and that $\Phi/N$ is bounded
above.  Constants in this Lemma may depend on these finite bounds.
With
$E_i=|\nabla N|^2+|\nabla^2f|^2$ and
$\rho=(1+y)/d$, one has
\begin{equation}
 \mathcal P:\nabla^2G\ge cE_i-C(1+|\dd G|^2).
 \label{eq:calculus-energy-nondivergence}
\end{equation}
There is a finite $k_0>0$, depending only on the stated bounds, such
that for $Y=s^{-k_0}$,
\begin{equation}
 \divg_g(\rho\mathcal P\dd Y)\ge cE_i-C.
 \label{eq:calculus-energy-divergence}
\end{equation}
These inequalities hold also at $\dd f=0$, and their constants use no
derivatives of $T$.

On a region where $p=0$, $s=N$, and $\Delta N=\Phi$, the corresponding
inequality is $\Delta Y\ge c|\nabla N|^2-C$.
At a constant-height gluing face satisfying $N_o=s_i$ and matching
flux of $B$, the directed inner flux $(\rho\mathcal P\dd Y)_n$ and
outer flux $\partial_nY$ differ by a bounded amount.  The bound
depends only on the finite field bounds and the fixed face geometry.
\end{lemma}

\begin{proof}
At nonzero gradient write $\gamma=\dd G$.  The identity
$\gamma=J\dd l-I\dd u$ gives
$\dd u=y\dd l-(1+y)\gamma$.
Inserting it into Equation~\eqref{eq:calculus-gradient-quadratic}
gives the exact expression
\begin{align}
 \mathcal Q={}&J\left[|U|^2+(1-y)l_1^2
                                  +(2-y)|\dd_\perp l|^2\right]\notag\\
 &-\frac{2y(2-y)}{1+y}\dd l\cdot\gamma
          +\frac{1+4y-y^2}{1+y}|\gamma|^2.
 \label{eq:calculus-energy-positive-form}
\end{align}
On the bounded field range, $1-y$ has a positive lower bound and
\[
 \frac J{m^2}=\frac{N^2(1-y)^2}{1+y},\qquad
 |\nabla^2f|^2=m^2(l_1^2+2|\dd_\perp l|^2+|U|^2).
\]
The first line of Equation~\eqref{eq:calculus-energy-positive-form}
therefore controls $|\nabla^2f|^2$.  Its mixed term is at most
$Cy|\dd l||\gamma|$ and is absorbed with an error $C|\gamma|^2$.
Moreover,
$|\dd u|^2\le C(y^2|\dd l|^2+|\gamma|^2)$, so after reducing the
positive coefficient the same lower bound controls $E_i$.

In Equation~\eqref{eq:calculus-master-gradient}, the trace term has
$|JC_T|\le C z$.  Its absolute value is bounded by
$C\sqrt{E_i}$, since $z$ is comparable to $m$ on this field range.
Equation~\eqref{eq:calculus-second-form-components}, or direct
differentiation of $p(N,\dd f)$, gives
$|\pi|\le C\sqrt{E_i}$.
The remaining terms in the master identity are bounded below by a
constant minus $C\sqrt{E_i}$.  Young's inequality proves
Equation~\eqref{eq:calculus-energy-nondivergence}.

There is no singularity at zero slope.  The expansion
$G=-\log N+\tfrac12N^2|\dd f|^2+O(|\dd f|^4)$ and the lapse equation
give, directly at a point with $\dd f=0$,
\[
 \mathcal P:\nabla^2G
 =N^2|\nabla^2f|^2+|\dd u|^2-\Phi/N
                    -N(\mathcal K+\ell g):\nabla^2f.
\]
This proves the same inequality there.  The matrices
$\mathcal P$ and $\rho\mathcal P$ are smooth functions of
$(N,\dd f)$, despite their expressions in terms of $e$; for example
$q-1=-2y/(1+y)$ cancels the denominator in $e\otimes e$.
Their derivatives consequently satisfy
$|\divg_g(\rho\mathcal P)|\le C\sqrt{E_i}$.

For $Y=e^{k_0G}$ the product rule yields
\begin{align*}
 \divg_g(\rho\mathcal P\dd Y)
 =k_0Y\bigl[\rho\mathcal P:\nabla^2G
      +\divg_g(\rho\mathcal P)\cdot\dd G
      +k_0\rho\mathcal P(\dd G,\dd G)\bigr].
\end{align*}
The mixed coefficient term is bounded by $C\sqrt{E_i}|\dd G|$.
Absorb half the positive $E_i$ term with Young's inequality, and then
choose $k_0$ so the final term absorbs the resulting multiple of
$|\dd G|^2$.  The upper and lower bounds on $s$ bound $Y$ on this
range and prove Equation~\eqref{eq:calculus-energy-divergence}.
Outside, direct differentiation gives
\[
 \Delta N^{-k_0}=k_0N^{-k_0}
          \left[(k_0+1)|\dd\log N|^2-\Phi/N\right],
\]
which proves the stated outer inequality.

On the face, $\rho q=d$ and
Equation~\eqref{eq:calculus-boundary-lapse} gives
\begin{equation}
 (\rho\mathcal P\dd Y)_n
   =-\frac{k_0YB_n}{s}
          +\frac{k_0Y}{d}(k_T\hat z-yH),
 \qquad
 (\partial_nY)_o=-\frac{k_0YB_n}{s}.
 \label{eq:calculus-energy-seam-flux}
\end{equation}
The second term in the inner expression is bounded under the stated
finite bounds; the first terms agree.  This proves the flux assertion.
\end{proof}

\begin{lemma}[Newton flux and freezing its coefficient]\label{lem:newton-flux}
In the setting of Lemma~\ref{lem:gradient-energy}, let
$\mathcal T_f=\nabla^2f-(\Delta f)g$.  The exact smooth vector identity
is
\begin{equation}
 \rho\mathcal P\nabla G
   =-\frac{\nabla N}{s}
      +\mathcal T_f\left(\frac{s^3}{N}\nabla f\right)
      +Ts\nabla f,
 \qquad
 \divg_g\mathcal T_f=\Ric_g(\nabla f,\cdot).
 \label{eq:calculus-newton}
\end{equation}
The coefficient of $\mathcal T_f$ is smooth at $\dd f=0$.

More specifically, in a fixed smooth face chart take compatible
reference fields
$\mathcal V_*=(n_*,a_*\dd t,s_*)$ for
$\mathcal V=(N_i,\dd f,N_o)$, where $t$ is inner signed unit normal
distance, $n_*>0$, and $s_*=s(n_*,|a_*|)$.
Assume the reference ranges and the actual fields lie in a fixed
bounded positive lapse and finite slope range.  Put
$Y_*=s_*^{-k_0}$ and
\[
 c_*=(k_0Y_*/s_*) ,\qquad
 \mathfrak b_*=k_0Y_*\frac{s_*^3}{n_*}a_*\nabla t.
\]
For the following assertion assume also that $\Phi$ is bounded
on each side.  Modify the inner and outer energy fluxes by
\begin{equation}
 \mathcal J_i=\rho\mathcal P\dd Y+c_*B-\mathcal T_f\mathfrak b_* ,
 \qquad \mathcal J_o=\nabla Y+c_*B.
 \label{eq:calculus-frozen-flux}
\end{equation}
They have bounded directed flux difference across the face, and
\begin{align}
 |\mathcal J|&\le C\bigl(1+|\mathcal V-\mathcal V_*|\sqrt E\bigr),
 \label{eq:calculus-frozen-flux-size}\\
 \divg\mathcal J&\ge c'E-C-C\delta_{\mathcal C}
 \label{eq:calculus-frozen-flux-divergence}
\end{align}
in the joined weak formulation, with $E=E_i$ inside and
$E=|\nabla N_o|^2$ outside.  The second line is understood with the
smooth metric densities in coordinates, and $\delta_{\mathcal C}$ is
surface measure on the face.  The same conclusion holds on an interior
chart with an arbitrary constant covector reference and no face term;
its metric-dependent reference coefficients are taken as smooth
functions of position.
\end{lemma}

\begin{proof}
At nonzero gradient the vector multiplying $\mathcal T_f$ can be
written as
\[
 \frac{y\nabla f}{d m^2}=
       \frac{s^3}{N}\nabla f,
 \qquad \frac{zT}{d}e=Ts\nabla f.
\]
For a transverse component, multiplication of the proposed identity
by $d$ gives $-u_a+yl_a$ on both sides.
For its first component, the two terms preceding $Ts\nabla f$ on the
right side, multiplied by $d$, are
\[
 -u_1-y\tr U=q(yl_1-u_1)-zT,
\]
by Equation~\eqref{eq:calculus-height-nondivergence}.
The last term adds $zT$, giving exactly the first component of the
left side.  This proves the first identity, including at zero gradient
by smoothness.  Commuting the covariant derivatives of $f$ gives
$\nabla^j\nabla^2_{ij}f=\nabla_i\Delta f+\Ric_{ik}\nabla^kf$,
which is the second identity.

Multiply the first identity by $k_0Y$.  In
Equation~\eqref{eq:calculus-frozen-flux}, the coefficients of
$\nabla N$ and of $\mathcal T_f$ then vanish at the reference state:
\begin{align*}
 \mathcal J_i={}&
   \left(c_*-\frac{k_0Y}{s}\right)\nabla N
   +\mathcal T_f\left(k_0Y\frac{s^3}{N}\nabla f-\mathfrak b_*\right)\\
 &+c_*(B-\nabla N)+k_0YTs\nabla f.
\end{align*}
The last line is bounded, since
$B-\nabla N=-(\mathcal K+\ell g)p$.
The two coefficient differences in the first line are Lipschitz in
the bounded field variables, so they are bounded by
$C|\mathcal V-\mathcal V_*|$.
On the outer side,
$\mathcal J_o=(c_*-k_0N_o^{-k_0-1})\nabla N_o$,
whose coefficient also vanishes at $N_o=s_*$.
This proves Equation~\eqref{eq:calculus-frozen-flux-size}.

The reference vector $\mathfrak b_*$ and its covariant derivative are uniformly
bounded in the fixed chart.  Equation~\eqref{eq:calculus-newton} gives
\[
 |\divg_g(\mathcal T_f\mathfrak b_*)|
 \le C(1+|\nabla^2f|).
\]
Adding $c_*B$ changes the divergence by the bounded scalar $c_*\Phi$.
Thus Equation~\eqref{eq:calculus-energy-divergence} and Young's
inequality give $\divg\mathcal J\ge c'E-C$ on each open side.
If reference coefficients vary smoothly with the metric in an interior
chart, their derivatives add only terms bounded by $C(1+\sqrt E)$;
these are absorbed in the same way.

Across the face, the energy flux difference is bounded by
Equation~\eqref{eq:calculus-energy-seam-flux}, and the added current has
matching normal flux.  The vector $\mathfrak b_*$ is normal to the face.  Its
Newton contraction there uses only
\[
 \mathcal T_f(n,n)=-\tr_{\mathcal C}\nabla^2f=-f_nH_i,
\]
which is bounded by the constant height trace and the slope bound.
Consequently the modified normal flux difference is bounded.  Integration
by parts on the two sides, whose face area densities agree, gives
the surface error in
Equation~\eqref{eq:calculus-frozen-flux-divergence}.
\end{proof}
